\chapter{Le cycle court : \texttt{forge-rapide}}
\source{skills/forge-rapide/SKILL.md}
Parcourir les étapes 0 à 6 pour ajouter un champ ou un filtre est disproportionné. Le cycle court garde la rigueur et allège la cérémonie.

\section{Quand l'utiliser}
Le cycle court est \textbf{refusé} si l'une des réponses suivantes est oui : nouveau bounded context ? nouvelle dépendance ? migration de données ? rupture du contrat OpenAPI ? plus d'une règle métier ? décision structurante (ADR) ? Dans ce cas, passez par \cmd{forge-spec} (cycle complet).

\section{Déroulement}
\begin{enumerate}
  \item \textbf{Mini-spec} : copiez \cmd{_TEMPLATE-spec-rapide.md} en \cmd{SPEC-NNN.md}. Elle contient une règle (avec sa source), un scénario Gherkin et l'\cmd{operationId} si l'API change. Taille \texttt{S}. Seul un humain la passe à \emph{Validée}.
  \item \textbf{Réalisation} : branche \cmd{feat/SPEC-NNN-slug}, contrat OpenAPI si l'API change, test d'abord, puis code.
  \item \textbf{Vérifications} : linter, tests, \cmd{check-glossary-terms.sh}, et \cmd{check-rapide-eligible.sh}.
\end{enumerate}

\section{Le garde-fou d'éligibilité}
Pour que le cycle court ne devienne pas une porte dérobée, \cmd{scripts/check-rapide-eligible.sh docs/02-specs/SPEC-NNN.md} échoue si :
\begin{itemize}
  \item la spec n'est pas de taille \texttt{S} ;
  \item plus de 8 fichiers sont modifiés par rapport à \cmd{main} ;
  \item un manifeste de dépendances change (\cmd{package.json}, \cmd{pyproject.toml}, \cmd{go.mod}\dots) ;
  \item une migration ou un fichier \cmd{.sql} apparaît.
\end{itemize}
En cas d'échec, le travail a dépassé le cadre : basculez vers le cycle complet.

\astuce{Le seuil de 8 fichiers est arbitraire et se règle à l'usage. Une modification de \cmd{package.json} est refusée même sans nouvelle dépendance : le script préfère le faux refus à l'oubli.}
